%%%%%%%%%%%%%%%%%%%%%%%%%%%%%%%%%%%%%%%%%%%%%%%%%%%%%%%%%%%%%%%%%%%%%%%%%%%%%%%
%  RobotSNAP -- paper skeleton
%
%  Target : Journal of Intelligent & Robotic Systems (Springer, journal 10846)
%           Topical Collection "Embodied Intelligence and Human-Centered
%           Adaptation in Autonomous Systems"  (collection hgfjccegje)
%           Submission deadline: 31 October 2026
%
%  Frame  : the collection invites "digital twin technologies for simulation
%           and validation". We claim pipeline and timing compatibility with a
%           real robot stack, NOT real-robot validation. Every claim below is
%           paired with the experiment that would falsify it (Table 4).
%
%  Build  : put the official Springer Nature LaTeX package in ./template/ so
%           that template/sn-jnl.cls resolves, then
%               pdflatex main
%               bibtex   main
%               pdflatex main
%               pdflatex main
%           (glossaries is loaded but unused; if you ever use it, add
%            \makeglossaries and run makeglossaries between the passes.)
%%%%%%%%%%%%%%%%%%%%%%%%%%%%%%%%%%%%%%%%%%%%%%%%%%%%%%%%%%%%%%%%%%%%%%%%%%%%%%%

\documentclass[pdflatex,sn-mathphys-num]{template/sn-jnl}% Math and Physical Sciences Numbered Reference Style
\usepackage{graphicx}%
\usepackage{multirow}%
\usepackage{amsmath,amssymb,amsfonts}%
\usepackage{booktabs}%
\usepackage{algorithm}%
\usepackage{algorithmicx}%
\usepackage{algpseudocode}%
\usepackage{xcolor}%
\usepackage{hyperref}%   (the base template declared it twice; once is enough)
\usepackage[xindy,acronyms,symbols]{glossaries}
\usepackage[counterclockwise, figuresright]{rotating}
\usepackage{tikz}
\usepackage{forest}
\usetikzlibrary{trees,positioning,shapes,shadows,arrows.meta}
\usetikzlibrary{arrows,patterns}
\usepackage{colortbl}

% --- skeleton helpers: delete before submission ------------------------------
\newcommand{\tbd}[1][TBD]{\textcolor{red}{#1}}
\newcommand{\ph}[1]{\textcolor{red}{\langle #1\rangle}}
\newcommand{\note}[1]{\textcolor{blue}{\footnotesize[#1]}}

% --- notation ---------------------------------------------------------------
\newcommand{\simstep}{\ensuremath{\Delta t}}
\newcommand{\ros}{\textsc{ros}\,2}

\begin{document}

\title[RobotSNAP: Robot Social Navigation Assessment Platform]{RobotSNAP: Robot Social Navigation Assessment Platform}
% Alternative titles that say "digital twin" and "validation" out loud, which
% is the wording the collection uses. Pick one before submission.
%   \title[RobotSNAP: A ROS 2 Digital Twin for Human-Centered Navigation Assessment]{...}
%   \title[Assessing Human-Centered Social Navigation in a Robot-Interface-Coupled Digital Twin]{...}

% The line below is copied verbatim from the base template. The official
% Springer Nature form is \author*[1]{...} (no backslash before the star); keep
% whichever your local sn-jnl.cls accepts.
\author\*[1]{\fnm{First} \sur{Author}}\email{author@example.com}
\affil\*[1]{\orgdiv{Department of Robotics}, \orgname{University Name}, \orgaddress{\city{City}, \country{Country}}}

\abstract{
Deciding whether a navigation policy is genuinely human-centered takes more
than a success rate. It takes scenarios that can be authored, replayed and
shared; human models whose fidelity is declared rather than assumed; metrics
whose sensitivity to that fidelity is known; and an interface a real robotic
stack can connect to without being rewritten. In practice these four
requirements are met by four different tools, and results are therefore hard to
compare and hard to reproduce. We present RobotSNAP, an assessment platform
that unifies them behind a single contract. RobotSNAP is a Unity digital twin
that exposes a fixed surface of ten topics carrying standard ROS message types,
published by a peer that is deliberately unaware of who is listening. The same
scenario file, the same seed and the same policy run either on a \ros{} graph
through \texttt{ros\_tcp\_endpoint}, or on a dependency-free Python transport
that speaks the identical wire protocol. We make three falsifiable claims.
(i) \emph{Cross-transport equivalence}: for a given scenario and seed, both
transports yield the same trajectories and the same metric values within
\ph{tolerance}. (ii) \emph{Stack reuse}: a third-party navigation stack
(\ph{name, version, source}) runs unmodified against the simulator, and the
measured control-loop budget of \ph{X}\,ms supports rates up to \ph{Y}\,Hz.
(iii) \emph{Human-model sensitivity}: replacing the crowd model with
\ph{dataset}-calibrated trajectories changes the ranking of \ph{K} policies
for \ph{metric}, so conclusions drawn from a synthetic crowd are not
transferable without declaring the human model. We report \ph{N} episodes over
\ph{M} scenario families and release the scenarios, the seeds, the containers
and the analysis scripts. RobotSNAP is positioned as a digital twin for
simulation and validation: kinematic and perceptual transfer to physical
hardware is out of scope and is stated as such.
}

\keywords{Human-aware navigation; Social robot navigation; Benchmarking and
reproducibility; Digital twins for robotics; ROS 2; Simulation validation;
Multi-agent simulation; Robot assessment platforms}

\maketitle

%%%%%%%%%%%%%%%%%%%%%%%%%%%%%%%%%%%%%%%%%%%%%%%%%%%%%%%%%%%%%%%%%%%%%%%%%%%%%%%
\section{Introduction}
\label{sec:intro}

% One paragraph of context, one paragraph of problem, one of gap, then the
% contribution list. Aim: a reader who stops after page 2 must know exactly
% what is claimed and what is not.

\subsection{Human-centered navigation needs evidence, not demonstrations}
\label{sec:intro:motive}
% Why social navigation is judged badly today: videos, cherry-picked runs,
% incomparable metric definitions, undeclared crowd models. Cite the metric
% critique line of work.

\subsection{The four requirements nobody meets at once}
\label{sec:intro:gap}
% 1. authorable + replayable scenarios (incl. real floor plans)
% 2. declared human-model fidelity
% 3. metrics with known sensitivity
% 4. an interface a real stack already speaks
% Point at Table 3 for where the existing platforms stop.

\subsection{Contributions}
\label{sec:intro:contrib}
\begin{itemize}
  \item \textbf{C1 --- A transport-neutral contract.} A surface of ten topics
        carrying standard message types, with per-robot namespacing, ROS
        frames and a published simulation clock, behind which the simulator
        cannot tell \texttt{ros\_tcp\_endpoint} from a dependency-free Python
        endpoint (\S\ref{sec:platform:contract}).
  \item \textbf{C2 --- Scenario and crowd authoring that scales to crowds.}
        Point and area goals, multi-waypoint routes, groups and formations,
        crowd modes that break formation, occupancy-grid and image import,
        deterministic seeding (\S\ref{sec:platform:authoring}).
  \item \textbf{C3 --- A measured equivalence between the two transports.}
        Not "we publish ROS topics" but a test: same seed, same policy, same
        metrics across the two bindings (\S\ref{sec:results:equivalence}).
  \item \textbf{C4 --- Reuse of an unmodified third-party stack, with its
        timing budget.} What a real pipeline needs the simulator to provide,
        and what it costs (\S\ref{sec:results:stack}).
  \item \textbf{C5 --- Evidence that crowd fidelity changes the verdict.}
        The result that makes this an assessment paper rather than a software
        report (\S\ref{sec:results:sensitivity}).
  \item \textbf{C6 --- Reproducibility artefacts.} Scenarios, seeds, analysis
        scripts and containers (\S\ref{sec:method:repro}).
\end{itemize}

\subsection{Scope and non-claims}
\label{sec:intro:scope}
% Explicit, early, and repeated in the discussion:
%   - no real-robot experiments; the claim is pipeline and timing
%     compatibility, verified against a stack that normally drives hardware;
%   - no human-subject study; the human side is validated against recorded
%     pedestrian trajectories, not against people in the loop;
%   - kinematic and perceptual transfer is not assessed.
% Editors forgive a declared limit; they do not forgive an untested claim.

\subsection{Organisation}
\label{sec:intro:org}

%%%%%%%%%%%%%%%%%%%%%%%%%%%%%%%%%%%%%%%%%%%%%%%%%%%%%%%%%%%%%%%%%%%%%%%%%%%%%%%
\section{Related Work}
\label{sec:related}

\subsection{Social navigation planners and human models}
\label{sec:related:planners}
% Reciprocal collision avoidance (ORCA), social force model, learning-based
% policies, group-aware planners. Keep it short: the paper is not about a
% planner, it is about how planners are judged.

\subsection{Benchmarks, simulators and platforms}
\label{sec:related:benchmarks}
% SocNavBench, SocialGym 1.0/2.0, SEAN 2.0, MRPB 1.0, Arena-Bench, Habitat 3.0,
% Jubileo, IBISCape. Then Table 3.

\subsection{Metric design and its critics}
\label{sec:related:metrics}
% The "human-aligned benchmarking" line: comfort, legibility, personal space;
% the argument that metric choice determines the verdict.

\subsection{Where RobotSNAP sits}
\label{sec:related:position}
% The gap: platforms either author scenarios well, or connect to real stacks
% well, rarely both; almost none test their own transport equivalence.

\begin{table*}[t]
  \centering
  \caption{Positioning against the closest platforms. Columns marked with a
  dash are either absent or not documented by the authors.
  \note{Confirm every cell against the primary source before submission; this
  is the table a reviewer will check first.}}
  \label{tab:positioning}
  \begin{tabular}{lcccccc}
    \toprule
    Platform & Scenario authoring & Real floor plans & Crowd / group modes
      & ROS \ros{} binding & Dependency-free Python & Transport equivalence tested \\
    \midrule
    SocNavBench        & \tbd & \tbd & \tbd & \tbd & \tbd & \tbd \\
    SocialGym 2.0      & \tbd & \tbd & \tbd & \tbd & \tbd & \tbd \\
    SEAN 2.0           & \tbd & \tbd & \tbd & \tbd & \tbd & \tbd \\
    MRPB 1.0           & \tbd & \tbd & \tbd & \tbd & \tbd & \tbd \\
    Arena-Bench        & \tbd & \tbd & \tbd & \tbd & \tbd & \tbd \\
    Habitat 3.0        & \tbd & \tbd & \tbd & \tbd & \tbd & \tbd \\
    Jubileo            & \tbd & \tbd & \tbd & \tbd & \tbd & \tbd \\
    \midrule
    \textbf{RobotSNAP} & yes & yes & yes & yes & yes & \textbf{yes} \\
    \bottomrule
  \end{tabular}
\end{table*}

%%%%%%%%%%%%%%%%%%%%%%%%%%%%%%%%%%%%%%%%%%%%%%%%%%%%%%%%%%%%%%%%%%%%%%%%%%%%%%%
\section{The RobotSNAP Platform}
\label{sec:platform}

\subsection{Design requirements from four usage profiles}
\label{sec:platform:requirements}
% The requirement set is derived from the intended users, not from taste.
% Table 1 is the contract with the reader: each profile lists what it must not
% be forced to install.

\begin{table}[t]
  \centering
  \caption{The four usage profiles the platform must serve, and the dependency
  each one is allowed to have. The fourth is the union of the first three and
  is the reason the contract must be transport-neutral rather than merely
  ROS-compatible.}
  \label{tab:profiles}
  \begin{tabular}{p{0.20\linewidth}p{0.26\linewidth}p{0.42\linewidth}}
    \toprule
    Profile & Entry point & Must not require \\
    \midrule
    Software / HRI prototyping & Dependency-free Python client & Any ROS installation \\
    Roboticist & \ros{} graph through \texttt{ros\_tcp\_endpoint} & Python client, Gymnasium \\
    Learning / RL researcher & Gymnasium environment, PyTorch & ROS, hand-written plumbing \\
    Full pipeline & All of the above, simultaneously & Rewriting a policy per transport \\
    \bottomrule
  \end{tabular}
\end{table}

\subsection{System architecture}
\label{sec:platform:architecture}
\begin{figure}[t]
  \centering
  % \includegraphics[width=\linewidth]{fig/architecture.pdf}
  \fbox{\parbox{0.92\linewidth}{\centering\vspace{1.5cm}
  \textbf{Figure 1} --- architecture.\\[2mm]
  Simulator core (scene, physics, sensors, scenario runtime) $\to$
  transport (identical wire protocol) $\to$ either \texttt{ros\_tcp\_endpoint}
  on a \ros{} graph, or \texttt{robotsnap.bridge} in pure Python.\\
  Both feed the same client: a policy, a Gymnasium environment, a recorder.
  \vspace{1.5cm}}}
  \caption{One scene, one contract, two bindings. The transport is the only
  component that changes between the top and the bottom of the stack.}
  \label{fig:architecture}
\end{figure}
% Layers: simulation core; scenario and map authoring; agent layer (robot(s),
% humans, groups); transport; client bindings. Keep to one page.

\subsection{The topic contract}
\label{sec:platform:contract}

\begin{table}[t]
  \centering
  \caption{The complete surface: ten topics, standard message types, no custom
  \texttt{.msg} file. Every per-robot stream is additionally reachable at
  \texttt{/robot\_\textless id\textgreater/\textless topic\textgreater}; the
  unprefixed name reaches the primary robot only. This is what makes the
  dependency-free Python path possible.}
  \label{tab:topics}
  \begin{tabular}{llll}
    \toprule
    Topic & Type & Dir. & Content \\
    \midrule
    \texttt{/clock}                        & \texttt{rosgraph\_msgs/Clock}    & out & simulation clock \\
    \texttt{/odom}                         & \texttt{nav\_msgs/Odometry}      & out & robot pose and twist \\
    \texttt{/scan}                         & \texttt{sensor\_msgs/LaserScan}  & out & lidar \\
    \texttt{/map}                          & \texttt{nav\_msgs/OccupancyGrid} & out & occupancy grid, walls 100 / free 0 \\
    \texttt{/cmd\_vel}                     & \texttt{geometry\_msgs/Twist}    & in  & velocity command \\
    \texttt{/simulation/state}             & \texttt{std\_msgs/String} (JSON) & out & whole-session snapshot \\
    \texttt{/simulation/agents}            & \texttt{std\_msgs/String} (JSON) & out & every agent, in the robot frame \\
    \texttt{/simulation/control}           & \texttt{std\_msgs/String} (JSON) & in  & session and crowd commands \\
    \texttt{/simulation/control\_result}   & \texttt{std\_msgs/String} (JSON) & out & answer to one command \\
    \texttt{/reset\_done}                  & \texttt{std\_msgs/Bool}          & out & handshake: world ready \\
    \bottomrule
  \end{tabular}
\end{table}
% Frames, units and conventions paragraph: ROS frame (x forward, y left, z up,
% yaw in radians); conversion happens inside the simulator so no client
% transforms anything; image order for the occupancy grid.

\subsection{Two bindings, one wire protocol}
\label{sec:platform:bindings}
% The simulator dials out as a client. Whatever it dials answers: either the
% ROS endpoint, or the pure-Python endpoint seeded from the same type table.
% Consequence: one scene runs in both setups and the simulator cannot tell the
% two apart. That property is a design claim here and a tested claim in
% \S\ref{sec:results:equivalence}.

\subsection{Scenario and map authoring}
\label{sec:platform:authoring}
% Occupancy-grid import (PNG and HouseExpo layouts); point goals and area
% goals; multi-waypoint routes for humans; groups and formations; crowd modes:
% synchronized, independent, and the per-agent scatter used for crossing and
% circle crowds; loop / wait / vanish terminal behaviours; per-route and
% per-scenario seeding. State plainly the rule that a goal is either a point
% or an area, never both.

\subsection{Human modelling}
\label{sec:platform:humans}
% Two modes only, and say so: an in-simulator social-force model, and external
% control from a client (this is the hook a \ros{} policy uses). Formation and
% conflict behaviour, group cohesion, and the deliberate absence of learned
% human models in this version.

\subsection{Heterogeneous multi-robot support}
\label{sec:platform:multirobot}
% Several robots, several platforms (differential-drive and omnidirectional
% archetypes), per-robot command topic, per-robot control authority: a robot
% moves only while it is the selected one and stops when no command arrives.

\subsection{Runtime introspection}
\label{sec:platform:introspection}
% What an assessor actually needs: a toggleable debug overlay, recording, a
% live occupancy minimap with occlusion-aware perception cones, and a camera
% that keeps visibility of the tracked agent behind walls.

%%%%%%%%%%%%%%%%%%%%%%%%%%%%%%%%%%%%%%%%%%%%%%%%%%%%%%%%%%%%%%%%%%%%%%%%%%%%%%%
\section{Evaluation Methodology}
\label{sec:method}

\subsection{Claims and their falsification criteria}
\label{sec:method:claims}
% This table is the spine of the paper. Write the paper so that a reader can
% read Table 4 alone and know what was tested.

\begin{table*}[t]
  \centering
  \caption{Each claim, the instrument that tests it, and the outcome that
  would falsify it. Nothing in the results section is reported that is not
  anchored here.}
  \label{tab:claims}
  \begin{tabular}{p{0.22\linewidth}p{0.32\linewidth}p{0.38\linewidth}}
    \toprule
    Claim & Test & Falsified if \\
    \midrule
    C3 Cross-transport equivalence
      & Same scenario, same seed, same policy, replayed through
        \texttt{ros\_tcp\_endpoint} and through the Python endpoint; compare
        full state streams and every metric
      & Any trajectory or metric differs beyond \ph{tolerance}, or the
        difference grows with episode length \\
    C4 Stack reuse
      & Run \ph{stack} unmodified, recording every source edit forced by the
        interface
      & The stack needs a change to its navigation behaviour, or the control
        loop misses its deadline at the stated rate \\
    C5 Human-model sensitivity
      & Rank \ph{K} policies under a synthetic crowd and under
        \ph{dataset}-calibrated trajectories
      & The ranking is identical across all metrics, i.e.\ crowd fidelity
        does not matter \\
    \bottomrule
  \end{tabular}
\end{table*}

\subsection{Scenario suite}
\label{sec:method:scenarios}
% Families chosen so that each isolates one social phenomenon. Say how many
% seeds per family and why.

\begin{table}[t]
  \centering
  \caption{Scenario families. The first five are authored from occupancy
  grids; IMT is an existing Unity scene, which is what makes the suite cover
  both worlds rather than only the synthetic one.}
  \label{tab:scenarios}
  \begin{tabular}{llll}
    \toprule
    Family & Agents & Phenomenon isolated & Seeds \\
    \midrule
    Straight corridor      & 1 + 2  & head-on passing            & \tbd \\
    Intersection           & 1 + 3  & priority negotiation       & \tbd \\
    Perpendicular crossing & 1 + 4  & crossing flow              & \tbd \\
    Crowd crossing         & 1 + 12 & dense anti-flow            & \tbd \\
    Circle crowd           & 1 + 10 & encircling group           & \tbd \\
    IMT (Unity scene)      & 1 + \ph{n} & real topology          & \tbd \\
    \bottomrule
  \end{tabular}
\end{table}

\subsection{Metrics}
\label{sec:method:metrics}
% Definitions matter more than names: a reviewer will check that each metric
% is computable from the recorded streams and that its units are stated.

\begin{table}[t]
  \centering
  \caption{Metrics, computed from the recorded topic streams so that a third
  party can recompute them from a published recording without re-running the
  simulator.}
  \label{tab:metrics}
  \begin{tabular}{llll}
    \toprule
    Metric & Definition & Unit & Aggregation \\
    \midrule
    Success rate        & reached final goal within the time budget & \%        & mean over seeds \\
    Collisions          & agent contact events                       & count     & sum, then mean \\
    Time to goal        & mission duration                           & s         & median, IQR \\
    Path efficiency     & shortest feasible path / travelled path    & ratio     & mean \\
    Min.\ separation    & minimum robot--human distance              & m         & min, then mean \\
    Personal-space cost & integral of violation of \ph{radius}       & \ph{unit} & mean \\
    Comfort             & \ph{definition: jerk, or acceleration}     & \ph{unit} & mean \\
    Frozen time         & time with speed below \ph{threshold}       & s         & mean \\
    \bottomrule
  \end{tabular}
\end{table}

\subsection{Statistical protocol}
\label{sec:method:stats}
% Paired comparisons over seeds, since the same seed makes scenarios
% comparable; report effect sizes and confidence intervals, not only p-values;
% state the correction used for the family of comparisons.

\subsection{Reproducibility artefacts}
\label{sec:method:repro}
% Containers, launch files, seeds, the analysis scripts, and the machine used
% to take the timing measurements.

%%%%%%%%%%%%%%%%%%%%%%%%%%%%%%%%%%%%%%%%%%%%%%%%%%%%%%%%%%%%%%%%%%%%%%%%%%%%%%%
\section{Results}
\label{sec:results}

\subsection{Cross-transport equivalence}
\label{sec:results:equivalence}
% The headline. Table: per metric, ROS-transport value, Python-transport
% value, absolute and relative difference. Then one figure: the divergence of
% the state streams over episode time, which is the plot that shows whether
% the difference is bounded or accumulating.

\subsection{An unmodified third-party stack}
\label{sec:results:stack}
% What was reused as-is, what had to be configured, and what the vendor would
% have had to change on hardware. Report the count of source edits: it should
% be zero in the navigation logic, and non-zero only in launch and parameter
% files. That distinction is the content of the claim.

\subsection{Control-loop budget}
\label{sec:results:timing}
\begin{table}[t]
  \centering
  \caption{Where a control step goes. The right column is the part a roboticist
  actually needs: which loop rates the platform can support.}
  \label{tab:timing}
  \begin{tabular}{lcc}
    \toprule
    Stage & Mean & 95th percentile \\
    \midrule
    Simulator step (\simstep{} = \ph{value}\,s)      & \tbd & \tbd \\
    Bridge round trip                                & \tbd & \tbd \\
    \ros{} publish $\to$ client receive              & \tbd & \tbd \\
    Client compute                                   & \tbd & \tbd \\
    End-to-end control period                        & \tbd & \tbd \\
    \midrule
    Sustained rate                                   & \ph{Y}\,Hz & --- \\
    \bottomrule
  \end{tabular}
\end{table}

\subsection{Does crowd fidelity change the verdict?}
\label{sec:results:sensitivity}
% The human-centered result. Rank policies under the synthetic crowd, then
% under recorded trajectories; report rank correlation and the cases where the
% order flips. One figure: rank before versus rank after, with the policies
% that swap labelled.

\subsection{Scenario-family results and the ranking of baselines}
\label{sec:results:families}
% The conventional results table, per family, per policy. It is here because
% it is expected, not because it is the contribution.

\subsection{Scaling}
\label{sec:results:scaling}
% Agents versus achievable step rate, on the stated machine. State the crowd
% size used for the timing measurements, and whether the timed episodes were
% in motion or had already converged --- a timing figure taken after the agents
% have arrived is not a timing figure.

%%%%%%%%%%%%%%%%%%%%%%%%%%%%%%%%%%%%%%%%%%%%%%%%%%%%%%%%%%%%%%%%%%%%%%%%%%%%%%%
\section{Discussion}
\label{sec:discussion}

\subsection{What the design decisions buy}
\label{sec:discussion:design}
% Connect each requirement of Table 1 to the design choice that satisfies it,
% and to the result that evidences it.

\subsection{Limits of the validation}
\label{sec:discussion:limits}
% State, in this order and without hedging:
%   - the compatibility demonstrated is interface and timing compatibility;
%   - kinematic transfer, sensor transfer and actuator transfer are not
%     assessed, and nothing here licenses a sim-to-real claim;
%   - the human side is validated against recorded trajectories, not against
%     people in the loop; a human-subject study is future work;
%   - one third-party stack is not a population of stacks.

\subsection{Threats to validity}
\label{sec:discussion:threats}
% Construct (do the metrics measure what the names say), internal (seed reuse,
% measurement conditions), external (topologies, crowd densities, robot
% platforms).

\subsection{Towards automated, statistical benchmarking}
\label{sec:discussion:roadmap}
% The direction this platform is built for: policy in, dataset of layouts
% (HouseExpo and successors), ladder of topology families, statistical report
% out. Say what is missing today: the multi-agent learning loop, and the real
% deployment that would close the last gap.

\subsection{Recommendations for reporting social navigation results}
\label{sec:discussion:reporting}
% Short, opinionated checklist: declare the human model, its calibration and
% its provenance; publish seeds; report to whom the metrics refer; state the
% transport and the loop rate at which the policy ran. A platform paper earns
% its place partly by leaving the field a reporting norm.

%%%%%%%%%%%%%%%%%%%%%%%%%%%%%%%%%%%%%%%%%%%%%%%%%%%%%%%%%%%%%%%%%%%%%%%%%%%%%%%
\section{Conclusion}
\label{sec:conclusion}
% One paragraph on what was built, one on what was tested, one on what is not
% claimed, one sentence on the next step.

%%%%%%%%%%%%%%%%%%%%%%%%%%%%%%%%%%%%%%%%%%%%%%%%%%%%%%%%%%%%%%%%%%%%%%%%%%%%%%%
\bmhead{Declarations}
% If \bmhead is not defined by your local class, use \section*{Declarations}.

\begin{itemize}
  \item Funding: \tbd
  \item Competing interests: \tbd
  \item Ethics approval: not applicable --- no human participants.
  \item Data availability: scenarios, seeds, analysis scripts and containers
        are released at \ph{URL}.
  \item Code availability: same repository; the simulator version used for
        every number in this paper is tagged \ph{tag}.
  \item Authors' contributions: \tbd
\end{itemize}

\bmhead{Appendix A --- Complete topic reference}
% Extend Table 2 with the per-robot naming rule, frame conventions, and the
% JSON schemas of the four simulation topics.

\bmhead{Appendix B --- Scenario file schema}
% One commented example, showing a point goal, an area goal, a multi-waypoint
% human route, a group, and a crowd mode.

\bmhead{Appendix C --- Metric definitions}
% Exact formulas, so a third party can recompute every number from a
% recording.

\bmhead{Appendix D --- Reproduction recipe}
% The shortest path from a clean machine to Table 5: commands, containers,
% expected runtime.

\bibliography{sn-bibliography}

\end{document}
